\clearpage

\section{JUSTIFICACIÓN}

    El presente proyecto surge de la necesidad de establecer las condiciones operativas y tecnológicas que permitan al Desafío Bebras desarrollarse de manera organizada y sostenible en Bolivia, considerando que su ejecución involucra procesos coordinados que abarcan desde la inscripción de unidades educativas escolares hasta la realización
    de las pruebas y la publicación de resultados dirigidos a estudiantes de nivel primario y secundario. La correcta articulación de estos procesos resulta determinante para garantizar una experiencia consistente para los estudiantes participantes y para asegurar la continuidad del concurso en ediciones posteriores. Las razones que motivan el desarrollo de la propuesta se fundamentan en dos dimensiones complementarias: la práctica y la metodológica.

    Desde el plano práctico, el proyecto ofrece una contribución directa a la organización del Desafío Bebras en Bolivia, al proponer una plataforma orientada a gestionar los procesos concretos que exige su ejecución anual en el ámbito escolar. Su aplicación permitirá coordinar con mayor eficiencia el registro de instituciones educativas de nivel primario y
    secundario, la administración de categorías por rango etario escolar, la realización de pruebas y el seguimiento de los estudiantes participantes. De esta manera, se busca fortalecer la capacidad organizativa nacional para operar el concurso de forma autónoma y estructurada, en correspondencia con los lineamientos establecidos por la comunidad internacional de Bebras. Dado que la plataforma será validada con usuarios reales vinculados al concurso, los resultados obtenidos constituirán evidencia concreta sobre su aporte a la viabilidad y continuidad del Desafío Bebras entre niños y jóvenes en edad escolar del contexto boliviano.

    Desde el plano metodológico, el proyecto se desarrollará aplicando la metodología ágil Extreme Programming (XP), adecuada para equipos pequeños de dos desarrolladores, lo que permitirá una construcción iterativa e incremental del sistema mediante iteraciones cortas y prácticas de programación en pareja e integración continua, facilitando la incorporación de retroalimentación por parte de los usuarios \cite{beckandres2004}. Para el análisis del sistema se emplearán técnicas de levantamiento de requerimientos que incluyen entrevistas con el equipo organizador nacional y revisión de los procesos establecidos por la comunidad internacional de Bebras. El diseño del sistema seguirá principios de arquitectura modular cliente-servidor, aplicando modelado de base de datos relacional y diseño de interfaces orientado a la usabilidad en contextos escolares. La validación se realizará mediante pruebas funcionales estructuradas, permitiendo contrastar objetivamente la contribución de la plataforma a la organización del concurso. En este sentido, la propuesta no solo busca materializar una solución tecnológica, sino también constituirse en una referencia metodológica para el desarrollo de
    plataformas educativas orientadas a competencias de pensamiento computacional en contextos escolares similares.